\documentclass[11pt,a4paper]{article}
\usepackage[utf8]{inputenc}
\usepackage{amsmath,amssymb}
\usepackage{fancyhdr}
\usepackage[margin=2.5cm]{geometry}

\pagestyle{fancy}
\fancyhf{}
\fancyhead[R]{Ivan Mišustin}
\renewcommand{\headrulewidth}{0pt}

\begin{document}

\begin{center}
    \LARGE\textbf{Lineární algebra -- domácí úkol č. 1}
\end{center}
\vspace{1em}

Pro $n \in \mathbb{N}$ určete mocninu matice:
\[
  \begin{pmatrix}
    1 & 1 & 0 \\
    -1 & 1 & 0 \\
    0 & 0 & 1
  \end{pmatrix}^n
\]

Pro libovolnou mocninu $n \in \mathbb{N}$ v matici zůstane poslední řádek a sloupec stejný, jelikož jejich nulové prvky jsou vždy výsledkem skalárního součinu ve formě:
\[
  \begin{pmatrix} 0 \\ 0 \\ 1 \end{pmatrix} \cdot \begin{pmatrix} a \\ b \\ 0 \end{pmatrix} = 0
\]
nebo
\[
  \begin{pmatrix} a \\ b \\ 0 \end{pmatrix} \cdot \begin{pmatrix} 0 \\ 0 \\ 1 \end{pmatrix} = 0
\]
a jednotkový prvek v pravém spodním rohu tedy vždy vzniká skalárním součinem ve formě:
\[
  \begin{pmatrix} 0 \\ 0 \\ 1 \end{pmatrix} \cdot \begin{pmatrix} 0 \\ 0 \\ 1 \end{pmatrix} = 1
\]

Z toho mohu tento problém zredukovat na následující jednodušší problém:
\[
  A = \begin{pmatrix}
    1 & 1 \\
    -1 & 1
  \end{pmatrix} \quad\text{a tedy hledám}\quad A^n = \begin{pmatrix}
    1 & 1 \\
    -1 & 1
  \end{pmatrix}^n
\]

Prvních pár mocnin této matice:
\begin{align*}
  A^2 &= \begin{pmatrix}
    0 & 2 \\
    -2 & 0
  \end{pmatrix}, \\
  A^3 &= \begin{pmatrix}
    -2 & 2 \\
    -2 & -2
  \end{pmatrix}, \\
  A^4 &= \begin{pmatrix}
    -4 & 0 \\
    0 & -4
  \end{pmatrix}.
\end{align*}

Vidím že $A^4$ je $(-4)$-násobkem identitní matice $\mathbb{E}_2$, a tedy každá mocnina která je násobkem $4$ se zapíše následovně:
\begin{align*}
  &\text{pro } n = 4q, \\
  &A^n = (-4)^q \mathbb{E}_2.
\end{align*}

Z toho tedy plyne že pro libovolnou mocninu $n$ mi záleží hlavně na zbytku po dělení $n$ čtyřmi, kde výsledná matice $A^n$ bude vypadat následovně:
\begin{align*}
  n &= 4q + r, \quad\text{kde } r \in \{0,1,2,3\}, \\
  A^n &= (-4)^q A^r.
\end{align*}

Nakonec pro původní matici:
\begin{align*}
  \begin{pmatrix}
    1 & 1 & 0 \\
    -1 & 1 & 0 \\
    0 & 0 & 1
  \end{pmatrix}^n
  &\quad\text{kde } n = 4q + r; \quad r \in \{0,1,2,3\} \\
  &= (-4)^q \begin{pmatrix}
    1 & 0 & 0 \\
    0 & 1 & 0 \\
    0 & 0 & (-4)^{-q}
  \end{pmatrix} \quad\text{pro } r = 0, \\[1ex]
  &= (-4)^q \begin{pmatrix}
    1 & 1 & 0 \\
    -1 & 1 & 0 \\
    0 & 0 & (-4)^{-q}
  \end{pmatrix} \quad\text{pro } r = 1, \\[1ex]
  &= (-4)^q \begin{pmatrix}
    0 & 2 & 0 \\
    -2 & 0 & 0 \\
    0 & 0 & (-4)^{-q}
  \end{pmatrix} \quad\text{pro } r = 2, \\[1ex]
  &= (-4)^q \begin{pmatrix}
    -2 & 2 & 0 \\
    -2 & -2 & 0 \\
    0 & 0 & (-4)^{-q}
  \end{pmatrix} \quad\text{pro } r = 3.
\end{align*}

\end{document}
